% ===========================================================================
\chapter{\label{chap:discussao}Discussão}

Este capítulo interpreta os resultados à luz da literatura do Capítulo~\ref{chap:fund}. Como a interpretação depende do que os dados mostrarem, cada seção explicita os mecanismos em jogo e o que cada desfecho plausível significaria; \resultado{após a coleta, manter em cada seção apenas a interpretação correspondente ao desfecho observado e reescrever os parágrafos condicionais como afirmações}. As Seções~\ref{sec:limitacoes} e~\ref{sec:etica} tratam das limitações e das questões éticas.

\section{A ancoragem funciona, e o que ela não resolve}

A comparação entre a proposta (P) e a linha de base sem fonte (S) é o teste central da primeira promessa. O mecanismo esperado é direto: com o fragmento no prompt, o modelo tem o conteúdo exato a ser perguntado, e sua tarefa se aproxima da reformulação; sem o fragmento, ele depende do que sabe sobre o tema, e a questão tende a cobrar conhecimento que o material do estudante pode não conter. Se P superar S sobretudo no critério de ancoragem, mas não no de correção, o resultado indicará que o ganho da ancoragem não é tornar as questões mais corretas em termos absolutos, e sim torná-las justas: perguntar o que o estudante efetivamente leu. Essa distinção importa no estudo autônomo, em que uma questão correta mas alheia ao material não é um erro do modelo, mas é uma falha pedagógica. \resultado{interpretação do desfecho observado para P vs.\ S}.

A ablação com vizinhos (P+V) testa uma intuição comum em sistemas RAG: a de que mais contexto ajuda. Há um mecanismo em sentido contrário: fragmentos vizinhos podem conter informação que o fragmento-fonte não tem, e o modelo pode construir a questão sobre ela, afastando-se do trecho que deveria ancorá-la. Se P+V reduzir a ancoragem em relação a P, o resultado sugerirá que, para geração de questões, a recuperação deve servir para escolher o que perguntar, e não para enriquecer a geração, o que contraria o uso típico descrito na literatura de RAG. \resultado{interpretação do desfecho observado para P vs.\ P+V}.

A comparação com o gerador por regras (R) situa o ganho trazido pelo modelo de linguagem em relação ao que a geração automática produzia antes dos LLMs. Espera-se que R produza questões ancoradas por construção, porque copia frases do fragmento, mas com distratores pouco plausíveis e nível cognitivo concentrado em lembrar. Se esse padrão se confirmar, ele ilustra por que ancoragem, sozinha, não é critério suficiente de qualidade, e por que a rubrica precisou ser conjuntiva. \resultado{interpretação do desfecho observado para P vs.\ R}.

Os defeitos típicos e as métricas automáticas indicam o que ainda não está resolvido. Um viés de extensão acima do esperado, por exemplo, significaria que parte dos acertos do estudante pode vir de uma pista formal e não do conteúdo, o que comprometeria a função de verificação das questões sem aparecer em nenhuma inspeção individual. Já a concordância entre o juiz automático e os avaliadores humanos responde a uma pergunta prática: se, em uso real, sem revisão humana, é possível filtrar automaticamente as questões não ancoradas antes de exibi-las. Uma concordância substancial tornaria esse filtro um próximo passo natural; uma concordância fraca reforçaria a cautela de Zheng et al.\ quanto ao uso de modelos como juízes. \resultado{interpretação do desfecho observado para as métricas automáticas}.

\section{Onde a execução no navegador se sustenta}

O resultado da QP2 deve ser lido como um mapa, não como uma média. O mecanismo que o explica é a diferença de custo entre os dois modelos: a detecção facial é leve e roda sobre um vídeo de baixa resolução, enquanto o YOLO, mesmo em sua menor variante, processa imagens de 640$\times$640 \estr{pixels} em WASM com uma única \estr{thread}. Espera-se, portanto, que a detecção facial seja viável em quase todas as configurações e que o YOLO se rompa primeiro nas máquinas de entrada. \resultado{interpretação do desfecho observado: onde cada modelo deixou de ser viável}.

Esse padrão, se confirmado, dá sustentação empírica à degradação graciosa como decisão de projeto. A alternativa, exigir um computador mínimo, excluiria justamente os estudantes com equipamentos mais modestos, e a degradação preserva para eles o indicador principal e as questões. O custo é uma assimetria de informação: estudantes com computadores mais fracos recebem um painel sem eventos de celular. Esse custo é aceitável para uma ferramenta de automonitoramento, mas seria um problema sério se os indicadores fossem usados para comparar estudantes entre si, o que reforça a decisão de não haver professor nem ranqueamento no sistema.

Duas melhorias técnicas decorrem diretamente da análise e podem deslocar a fronteira de viabilidade: reduzir a resolução de entrada do YOLO, trocando acurácia por latência, e executá-lo com aceleração gráfica (WebGPU ou WebGL), quando disponível. Ambas são trabalhos futuros mensuráveis com o mesmo protocolo.

\section{O que o \fs mede, e o que não mede}

A QP3 verifica se o indicador responde aos estados que foi desenhado para capturar. Dois segmentos do roteiro concentram as ambiguidades esperadas. O primeiro é o olhar para baixo (S4): o indicador o classifica como desatento, mas no estudo real olhar para baixo pode significar consultar anotações, que é estudo. O segundo é o celular sobre a mesa com o estudante atento (S8): o detector conta um evento porque reconhece o objeto, não o comportamento. Nenhum desses casos é um erro de implementação; ambos são limites do que o sinal visual pode distinguir, e o roteiro foi desenhado para torná-los visíveis. \resultado{interpretação do desempenho observado em S4 e S8}.

O limite mais importante, porém, não aparece no roteiro, porque não pode ser encenado de forma verificável: a divagação mental com os olhos na tela. Um \fs alto é compatível com atenção plena e com uma mente que abandonou o texto. Por isso, este trabalho adota o termo engajamento visual e evita falar em atenção. Essa posição reconhece que a câmera observa comportamento, e que a relação entre comportamento visível e processamento cognitivo é indireta. Ela também explica por que as questões são parte indispensável da ferramenta, e não um complemento: as respostas oferecem o sinal que a câmera não oferece, sobre o processamento do conteúdo. O cruzamento entre os dois sinais no painel é, portanto, a forma de compensar, em parte, a limitação de cada um.

A análise de sensibilidade responde a uma crítica possível ao indicador: a de que seus parâmetros foram escolhidos por tentativa. Se o desempenho variar pouco em torno dos valores padrão, as conclusões não dependem de um ajuste fino; se variar muito, os parâmetros precisam ser recalibrados em populações mais diversas antes de qualquer uso. A comparação entre limiar calibrado e fixo testa, em particular, a hipótese de projeto de que o valor da EAR com o olho aberto varia entre condições de gravação o bastante para exigir calibração individual. \resultado{interpretação do desfecho observado na análise de sensibilidade}.

\section{Privacidade como propriedade, e seus limites}

A QP4 muda o tipo de afirmação que o trabalho pode fazer sobre privacidade. Em vez de declarar que as imagens não saem do navegador, o trabalho oferece um teste que qualquer pessoa pode executar e que demonstrou ser capaz de detectar o vazamento quando ele existe. Essa é a leitura operacional do \estr{privacy by design} de Cavoukian: a proteção embutida na arquitetura torna-se verificável justamente por não depender de uma promessa. \resultado{interpretação do desfecho observado}.

O mesmo teste delimita o alcance dessa garantia. Ele verifica o tráfego da versão testada, não de versões futuras; por isso, deve ser incorporado à integração contínua do projeto e executado a cada alteração. Ele não cobre a segunda fronteira de privacidade, a do texto do material, que é enviado ao provedor do LLM por necessidade. E ele documenta um fluxo que não envolve imagens, mas não é neutro: o download dos arquivos do MediaPipe a partir de uma rede de distribuição de terceiros. Hospedar esses arquivos no próprio servidor da aplicação eliminaria esse contato e é uma correção simples.

\section{Implicações para o projeto de ferramentas de estudo}

Três implicações ultrapassam o \ferramenta. A primeira é que a ancoragem em geração de questões deve ser medida, e não presumida, e que uma linha de base com o mesmo modelo sem o fragmento é a comparação que isola seu efeito. A segunda é que sensoriamento no cliente é uma forma concreta de conciliar monitoramento e privacidade, desde que o projeto trate a heterogeneidade do \estr{hardware} como requisito, com degradação graciosa e medição explícita. A terceira é metodológica: vídeos roteirizados, páginas de \estr{benchmark} e testes de tráfego permitem avaliar propriedades essenciais de ferramentas educacionais antes de envolver estudantes, o que reduz o risco e o custo dos estudos com participantes que vierem depois.

\section{\label{sec:limitacoes}Limitações}

As limitações a seguir são específicas deste trabalho e indicam, cada uma, que tipo de conclusão os resultados não autorizam.

\begin{itemize}
  \item \textbf{Nenhuma medida de aprendizagem.} O trabalho não mede se o \ferramenta melhora a compreensão, a retenção ou a atenção. As QPs estabelecem propriedades necessárias, não suficientes, para esse efeito.
  \item \textbf{Corpus pequeno e escolhido.} Três materiais abertos, com 30 fragmentos cada, não representam a variedade dos PDFs que estudantes reais enviam, como apresentações digitalizadas, listas de exercícios ou textos em outros idiomas. PDFs digitalizados sem camada de texto não são sequer processados pela extração atual.
  \item \textbf{Avaliadores próximos ao projeto.} Os dois avaliadores das questões são o autor e o orientador. A inspeção cega à condição reduz, mas não elimina, o viés de quem conhece o sistema. A concordância reportada é entre esses dois avaliadores e pode não se repetir com outros.
  \item \textbf{Um único modelo gerador.} Os resultados da QP1 valem para o modelo e a versão avaliados, com o prompt descrito. Modelos de linguagem mudam com frequência, e o protocolo, mais do que os números, é o que se transfere.
  \item \textbf{Um único rosto no vídeo.} As gravações da QP3 mostram apenas o autor. Tom de pele, formato do rosto, óculos de outros tipos, barba, maquiagem e deficiências visuais ou motoras podem alterar o desempenho da detecção facial, e nada nos resultados permite afirmar que o indicador funciona igualmente para outras pessoas.
  \item \textbf{Comportamentos encenados.} Os estados do roteiro são representados deliberadamente e com transições nítidas. Comportamentos espontâneos são mais graduais e ambíguos, e o desempenho em uso real tende a ser inferior.
  \item \textbf{Amostra de \estr{hardware} por conveniência.} As configurações da QP2 não representam o parque de computadores dos estudantes. Os resultados descrevem onde a viabilidade se rompe, não com que frequência isso ocorre na população.
  \item \textbf{Aproximações geométricas.} As razões usadas como \estr{yaw} e \estr{pitch} não são ângulos e são sensíveis à posição da câmera em relação à tela; câmeras externas, posicionadas ao lado do monitor, podem deslocar a referência da calibração.
  \item \textbf{Privacidade verificada em uma versão.} O teste da QP4 cobre a versão avaliada e o fluxo de uma sessão típica, não todos os caminhos possíveis da aplicação.
\end{itemize}

\section{\label{sec:etica}Ética e transparência}

\textbf{Aprovação ética.} A avaliação proposta não envolve participantes humanos na condição de sujeitos de pesquisa: as questões são avaliadas pelo autor e pelo orientador, as gravações mostram apenas o autor, e as execuções de \estr{benchmark} coletam dados sobre máquinas, não sobre pessoas, a partir de um vídeo pré-gravado. Por essa razão, entende-se que o protocolo não exige submissão ao Comitê de Ética em Pesquisa. \definir{confirmar esse entendimento com o orientador e, se necessário, com o CEP-PUCRS, em especial quanto às execuções de \estr{benchmark} por colegas e à inspeção heurística}. O estudo com participantes descrito no Apêndice~\ref{ap:usuarios} exigirá aprovação prévia.

\textbf{Dados pessoais.} Na ferramenta, a imagem da câmera é processada apenas no navegador e só é ativada com a ação explícita do estudante. O servidor armazena as respostas e a telemetria agregada, associadas a sessões e não a identidades civis. O conteúdo dos PDFs enviados é transmitido, em fragmentos, ao provedor do modelo de linguagem, e o estudante deve ser informado disso antes do envio, inclusive para que não envie documentos com dados pessoais de terceiros.

\textbf{Riscos sociais.} Um indicador de engajamento pode ser reaproveitado para vigilância, por exemplo se uma instituição passasse a exigir o \fs como prova de estudo. O projeto do \ferramenta reduz esse risco por decisões de arquitetura: não há papel de professor ou administrador, não há comparação entre estudantes, e as imagens não saem do dispositivo. Ainda assim, o código é aberto, e esses limites não impedem que terceiros o modifiquem. O indicador também pode produzir erros sistemáticos para grupos específicos (Seção~\ref{sec:limitacoes}), o que o torna inadequado para qualquer uso avaliativo ou disciplinar.

\textbf{Conflitos de interesse e financiamento.} \definir{declarar ausência de conflitos de interesse e de financiamento, ou descrevê-los; informar se o uso da API do modelo de linguagem foi custeado pelo autor ou por algum programa}.

\textbf{Disponibilidade.} O código, o corpus, as questões com suas classificações, os registros de desempenho, os registros numéricos por quadro e os \estr{scripts} de análise são disponibilizados no repositório indicado na Seção~\ref{sec:reprod}. As gravações de vídeo não são publicadas.

\textbf{Uso de ferramentas de IA na escrita.} \definir{declarar se e como ferramentas de IA generativa foram usadas na redação ou revisão do texto e do código, conforme a orientação da PUCRS}.
